\chapter{El verificador de pruebas}
\label{cap:verificador}

Hace falta una fórmula del lenguaje objeto que diga «\(x\) es demostrable». La estrategia de Gödel
es indirecta y es la única que funciona: no se define la demostrabilidad, se define **verificar una
demostración** —que es una comprobación finita— y luego se dice «existe algo que pasa la
verificación».

\section{Una demostración es una lista de líneas}

\begin{matematica}[La forma de una prueba]
Una demostración es una lista \(p\) de líneas. Cada línea lleva su conclusión y la etiqueta de la
regla que la justifica, más las referencias a sus premisas. La cadena es válida si \emph{cada} línea
está bien formada y \emph{cada} premisa que invoca ya estaba probada antes.
\end{matematica}

Eso se escribe con cuatro símbolos del lenguaje objeto, y los cuatro son \textbf{opacos}:

\leanfrag{lineWF}
\leanfrag{premsOf}
\leanfrag{runFn}
\leanfrag{chainOk}

\noindent \(\obj{\mathrm{lineWF}}(\ell)\) dice que la línea \(\ell\) es una instancia legítima de
alguna regla; \(\obj{\mathrm{premsOf}}(\ell)\) da la lista de premisas que invoca;
\(\obj{\mathrm{runFn}}(c,p)\) acumula las conclusiones de la cadena \(p\) sobre lo ya verificado
\(c\); y \(\obj{\mathrm{chainOk}}(c,p)\) es la comprobación completa.

\begin{leccion}[De dónde salen los 107]
El capítulo~\ref{cap:qpp} contaba 34 axiomas de aritmética y \textbf{107 de maquinaria de
codificación}. Aquí es donde están: cada uno de estos símbolos opacos necesita sus ecuaciones, y las
reglas del cálculo son \textbf{21 etiquetas} que \(\obj{\mathrm{lineWF}}\) tiene que reconocer una
por una. Que la maquinaria triplique a la aritmética no es descuido: es lo que cuesta escribir un
verificador \emph{dentro} de la teoría que verifica.
\end{leccion}

\section{El predicado de demostrabilidad}

Con el verificador, la demostrabilidad es una sola fórmula, y es \(\Sigma_1\):

\leanfrag{provFormulaC}

\begin{matematica}[$\Prov$, por fin en el lenguaje objeto]
\[
  \Prov(x) \;:=\; \exists p.\ \obj{\mathrm{chainOk}}(\mathrm{nil}, p)
  \;\wedge\; x \obj{\in} \obj{\mathrm{runFn}}(\mathrm{nil}, p)
\]
\selee{existe una cadena de prueba $p$ que pasa la verificación desde cero, y entre las conclusiones
que acumula está $x$.}
\end{matematica}

Y aplicándolo al código de una fórmula concreta:

\leanfrag{provCodeC}

\begin{leccion}[Por qué $\Sigma_1$ y por qué importa]
La forma es «\(\exists p\) tal que \emph{algo decidible} sobre \(p\)»: eso es
\(\Sigma_1\). Y es la forma que hace funcionar la incompletitud, porque una teoría aritmética
razonable demuestra todos los \(\Sigma_1\) verdaderos — comprobar un testigo concreto es cálculo,
no razonamiento. Esa propiedad es la condición \textbf{D1}, y es el capítulo siguiente. Su versión
\emph{provable} —que la teoría demuestre esa capacidad \emph{de sí misma}— es \textbf{D3}, y es el
frente abierto del proyecto.
\end{leccion}

\begin{muro}[La asimetría del verificador]
El verificador es \emph{sólido} por construcción: si una cadena pasa, hay demostración. Lo que no es
gratis es el recíproco negativo — que si algo \emph{no} es demostrable, la teoría pueda refutar
todas las cadenas. Eso es lo que falta para la mitad $\noderives\lnot G$ del primer teorema, y tiene
capítulo propio.
\end{muro}
